\chapter{INTRODUCTION}

\section{Background}

In the upstream oil and gas exploration and production industry, subsurface borehole drilling and reservoir evaluation represent operations of immense technical complexity and high financial expenditure \citep{ref3, ref9}. Drilling a single deep exploratory or offshore development well frequently requires tens of millions of dollars in capital investment; consequently, asset exploration teams urgently require rapid, continuous diagnostic insight into the nature of penetrated formation fluids as drilling advances \citep{ref3}. Globally, rotary drilling is conducted on a continuous industrial scale, with thousands of active rigs operating simultaneously across international hydrocarbon basins \citep{ref9, ref14, bakerhughes2024}. Because drilling is an uninterrupted, high-cost operation, Gas While Drilling (GWD) analysis is deployed almost universally as the primary real-time formation evaluation technique \citep{ref9, ref12}. 

As the drill bit crushes the subsurface rock matrix, trapped formation gases are liberated into the circulating drilling mud and carried up to the surface. A specialized surface extraction trap and gas chromatograph continuously measure the volumetric concentrations (in parts per million or volumetric percentages) of light to intermediate alkane hydrocarbon gas fractions, specifically methane ($C_1$), ethane ($C_2$), propane ($C_3$), iso-butane ($iC_4$), normal-butane ($nC_4$), iso-pentane ($iC_5$), and normal-pentane ($nC_5$), alongside Total Gas ($TG$) \citep{ref12, haworth1985}. This chromatographic data provides multidisciplinary exploration teams, encompassing \textbf{well-site geologists}, \textbf{operations geologists}, and \textbf{petrophysicists} with the earliest diagnostic signature of subsurface hydrocarbon presence during the very first bit run \citep{ref9, ref12}.

Despite the continuous acquisition of GWD chromatographic data during the initial drilling pass, standard field evaluation workflows remain fundamentally hindered by a major operational bottleneck \citep{ref9, ref11}. In conventional field practice, asset teams treat GWD data primarily as a secondary safety and monitoring record. Definitive hydrocarbon fluid typing, distinguishing dry gas caps, volatile oil columns, heavy oil pay, and non-productive water sands is routinely delayed until post-drilling Wireline Logging (WL) or complex Logging While Drilling (LWD) suites can be deployed \citep{ref8, ref12}. However, conducting wireline logging requires pulling the entire drill string out of the borehole (tripping), which introduces days of non-productive rig standby time, severe borehole collapse hazards, differential sticking risks, and hundreds of thousands to millions of dollars in specialized logging contractor fees \citep{ref3, ref9}. If hydrocarbon fluid regimes could be reliably and deterministically forecasted directly from first-drilling GWD data, wellsite and operations teams could make immediate geosteering corrections, optimize casing seat depths, or terminate non-productive wells early before incurring massive wireline logging expenditures \citep{ref3, ref8}.

A primary reason GWD data has historically been underutilized for early fluid forecasting lies in the architectural and operational limitations of existing formation evaluation software \citep{ref8, wood2020}. Commercial platforms deployed across the industry, such as Schlumberger (SLB) Techlog, AspenTech Geolog, Advanced Logic Technology WellCAD, and RockWare LogPlot as well as generic open-source Python scripts operate predominantly as \textit{passive data plotters} \citep{wood2020}. Drafting applications like WellCAD and LogPlot excel at rendering geotechnical symbols and composite strip tracks, but contain zero built-in computational intelligence to derive multi-ratio petrophysical matrices or forecast fluid types. Meanwhile, comprehensive workstation suites like Techlog and Geolog treat mud log gas curves as simple visual tracks; executing multi-ratio petrophysical forecasting requires users to write complex routines in proprietary scripting consoles (e.g., Techlog Python Script Editor or Geolog Loglan macros) or develop compiled C++ DLL plugins \citep{wood2020}. Furthermore, these enterprise platforms are notoriously monolithic, heavyweight, and slow to configure, burdening asset teams with steep learning curves, cluttered multi-level menus, and cost-prohibitive annual licensing barriers (often exceeding \$20,000 to \$50,000 per seat) that hinder fast-paced field operations \citep{ref8, wood2020}.

In addition to software latency and complexity, formation evaluation workflows are governed by strict compliance with petroleum Standard Operating Procedures (SOP) and international reserve certification standards. While modern machine learning (ML) and deep neural network models have been applied to well log prediction, their non-transparent "black-box" nature obscures internal reasoning, rendering their predictions non-auditable and indefensible during formal reserve certification audits under the \citet{ref3} Petroleum Resources Management System (PRMS) \citep{ref1, ref3}. Consequently, subsurface asset teams face an unsatisfactory dilemma between non-auditable AI models and cumbersome workstation plotting software that cannot automate real-time GWD fluid forecasting.

This situation reveals a critical technological and operational gap: there is currently an absence of a dedicated, lightweight decision support platform that bridges raw first-drilling GWD data with automated multi-ratio petrophysical derivation and deterministic hydrocarbon fluid forecasting \citep{ref8, wood2020}. To overcome these challenges, this research develops a dedicated, open-source web-based decision support system using Python (Dash and Plotly). The proposed application simplifies the operational workflow, executing automated multi-ratio petrophysical calculations and deterministic fluid forecasting in seconds without proprietary software dependencies \citep{ref3, ref8, ref11, ref12}.

\section{Product Formulation}

The product developed in this research is an open-source, web-based deterministic decision support system engineered to transform raw, first-drilling Gas While Drilling (GWD) chromatographic data into real-time multi-ratio petrophysical intelligence and automated hydrocarbon fluid zone forecasts. The product formulation is explicitly defined through the following declarative statements:

\begin{enumerate}
    \item The product is a lightweight, web application that ingests raw mud logging chromatographic gas records ($C_1$ through $nC_5$ and $TG$) across standard tabular and industry formats (\texttt{.csv}, \texttt{.txt}, \texttt{.xlsx}, and \texttt{.las}), automatically isolating numerical data matrices from unstructured rig comments and company headers.
    
    \item The product incorporates an automated calculation engine that derives a comprehensive matrix of 16 petrophysical indicators including classic Pixler ratios ($R_1$--$R_5$), expanded Haworth show ratios (Wetness $W_h$, Balance $B_h$, Character $C_h$), butane multipliers ($\text{Ratio } iC_4, \text{Ratio } nC_4$), Dryness Ratio ($DR$), Carbon-Weighted Density Index ($I_{\text{carbon}}$), and composite screening indices ($GOW, GOW_{\text{noTG}}, WBS, GOR$) at every continuous depth increment.
    
    \item The product executes a deterministic, rule-based majority-voting classification algorithm that categorizes subsurface formations into discrete reservoir fluid regimes (\textbf{Gas Pay Zone}, \textbf{Oil Pay Zone}, or \textbf{Non-Bearing Sand}) with 100\% mathematical traceability, ensuring full compliance with SPE PRMS Chapter 4 reserve certification audit standards.
    
    \item The product provides an interactive, depth-synchronized 23-track well log visualization environment built with Plotly WebGL, adhering to API RP 31A and SPWLA strip log conventions, fitted dynamically to 100\% screen width with zero horizontal scrolling.
    
    \item The product incorporates interactive pipeline management modules the *Column Configuration and Display Manager* and the *Petrophysical Formula and Threshold Manager* allowing geoscientists to dynamically calibrate mathematical expressions, adjust localized cutoff thresholds, and customize track layouts in real time without modifying backend source code.
    
    \item The product delivers sub-5-second end-to-end execution latency across deep well trajectories ($> 3000\text{ m}$) and is deployed globally via a serverless cloud architecture on Vercel to eliminate commercial licensing barriers and local installation friction.
\end{enumerate}

\section{Product Scope and Limitations}

To ensure development feasibility, clear research boundaries, and targeted operational utility, the scope and limitations of the developed product are established as follows:

\begin{enumerate}
    \item \textbf{Input Data and Format Scope:} The software supports chromatographic gas data consisting of light alkanes ($C_1, C_2, C_3, iC_4, nC_4, iC_5, nC_5$) and Total Gas ($TG$) indexed against continuous Measured Depth ($MD$). File ingestion is structured for delimited ASCII files (\texttt{.csv}, \texttt{.txt}), Microsoft Excel workbooks (\texttt{.xlsx}), and Log ASCII Standard files (\texttt{.las} versions 1.2 and 2.0). Highly specialized binary formats (e.g., proprietary DLIS/LIS wireline archives) are outside the current ingestion scope.
    
    \item \textbf{Core Computational Logic:} The calculation engine strictly implements analytical, closed-form petrophysical formulas (Pixler, Haworth, Whittaker, and SPE GOR guidelines) and a deterministic majority-voting classifier. The system deliberately excludes stochastic, uninterpretable machine learning weights to maintain full auditability under SPE PRMS Chapter 4 standards.
    
    \item \textbf{Technology Stack and Environment:} The application backend and computation engine are built in Python 3 utilizing vectorized NumPy and Pandas libraries. The reactive frontend is constructed using the Dash framework, Dash Bootstrap Components (\texttt{dbc}), and Plotly WebGL. The system is deployed serverless on Vercel and runs client-side in standard modern web browsers (Google Chrome, Mozilla Firefox, Microsoft Edge, Safari) on desktop and ruggedized field workstations.
    
    \item \textbf{Target User Group:} The software is tailored specifically for operational subsurface professionals, including well-site geologists, operations geologists, logging engineers, petrophysicists, and petroleum engineering students.
    
    \item \textbf{Evaluation and Verification Scope:} System verification encompasses automated unit tests with synthetic benchmark matrices, computational latency profiling across well profiles up to 10,000 intervals ($15,000\text{ m}$), empirical classification validation across 16 verified formation intervals benchmarked against Wireline Formation Tester (MDT/RFT) records under strict commercial confidentiality protocols, and usability evaluation via the standardized 10-item System Usability Scale (SUS) with a practicing subsurface domain professional.
\end{enumerate}

\section{Research Objectives}

The primary objective of this research is to design, implement, and evaluate an open-source, web-based deterministic decision support application capable of transforming first-drilling Gas While Drilling (GWD) chromatographic data into real-time multi-ratio petrophysical indicators and automated hydrocarbon fluid zone forecasts, resolving the operational delays of wireline logging, the audit risks of black-box models, and the sluggish complexity of passive log plotting software.

To systematically achieve this primary objective, the specific research goals are formulated as follows:

\begin{enumerate}
    \item To develop an automated Python calculation engine that ingests chromatographic gas fractions ($C_1$ through $nC_5$) and derives 16 petrophysical indicators (Pixler ratios, Haworth show ratios $W_h, B_h, C_h$, butane multipliers, Dryness, Carbon Index, $GOW, WBS, GOR$), executing deterministic majority-voting logic to classify fluid regimes (Gas, Oil, Non-Bearing) at every depth increment.
    
    \item To construct an interactive, depth-synchronized 23-track web visualization interface using Dash and Plotly tailored for wellsite geologists, operations geologists, and petrophysicists, featuring dynamic pipeline controls (Column Manager and Formula/Threshold Manager) for live operational calibration.
    
    \item To establish documented compliance with petroleum Standard Operating Procedures (SOP) and international reserve certification standards by:
    \begin{itemize}
        \item Grounding calculation rules in ratified SPE and AAPG literature as immutable default baselines, paired with dynamic user-override capabilities.
        \item Fulfilling SPE/WPC/AAPG/SPEE/SEG PRMS Chapter 4 deterministic auditability mandates through 100\% mathematically traceable closed-form logic.
        \item Structuring multi-track log presentations in accordance with API RP 31A and SPWLA wellsite mud logging reporting standards.
    \end{itemize}
    
    \item To optimize end-to-end execution latency across raw file ingestion, 16-parameter calculation, and 23-track rendering to complete in under 5.0 seconds for deep well trajectories ($> 3000\text{ m}$), and validate system usability with practicing subsurface professionals to achieve a System Usability Scale score of $S_{\text{SUS}} \ge 68.0/100$.
\end{enumerate}

\section{Research Significance and Benefits}

The development of this dedicated decision support system provides tangible contributions across industrial, operational, and academic domains:

\subsection{Industrial and Operational Contributions}

\subsubsection{Early First-Drilling Fluid Intelligence}~\\
By forecasting hydrocarbon fluid contacts directly from GWD data during the first bit run, \textbf{well-site geologists}, \textbf{operations geologists}, and \textbf{petrophysicists} gain critical formation insights days before wireline logging tools can be mobilized. This enables proactive geosteering adjustments, optimized casing seat placement, and early well termination in non-productive structures, saving millions of dollars in rig standby and drilling costs.

\subsubsection{Workflow Streamlining and Speed Optimization}~\\
By replacing monolithic, slow-to-configure workstation software with an intuitive, single-click web dashboard, asset teams eliminate hours of manual script configuration and multi-step data preparation. Sub-5-second execution latency provides immediate decision support within tight operational drilling windows.

\subsubsection{Full Compliance with Industry SOP and PRMS Audit Standards}~\\
The software guarantees complete alignment with petroleum industry standards through:
\begin{itemize}
    \item \textit{Ratified Engineering Lineage:} Mathematical formulas strictly execute equations published in standard petroleum literature (\citet{haworth1985} for show ratios, \citet{ref13} for Pixler boundaries, \citet{whittaker1991} for continuous composite gas indices, and \citet{spe2020} for GOR classifications).
    \item \textit{Deterministic Auditability under SPE PRMS:} In full accordance with \citet{ref2} and \citet{ref3} PRMS Chapter 4 standards, every calculation and classification is 100\% traceable from raw input to final output, eliminating the audit failure risks of black-box machine learning models.
    \item \textit{Standardized Mud Log Rendering:} Output tracks adhere directly to API RP 31A and SPWLA strip log conventions for direct operational reporting.
\end{itemize}

\subsubsection{Elimination of Commercial Software Licensing Barriers}~\\
By delivering a free, open-source web application tailored specifically for automated GWD multi-ratio forecasting, this platform removes the requirement for independent operators, service companies, and consultants to purchase expensive workstation licenses (frequently exceeding \$20,000 to \$50,000 annually per seat) merely to perform mud log analytics \citep{wood2020}.

\subsection{Academic and Educational Contributions}

\subsubsection{Democratization of Advanced Petrophysical Computing}~\\
Commercial software licensing costs frequently exclude academic institutions and student researchers from advanced log evaluation tools. Publishing this application under an open-source license democratizes access to multi-parameter petrophysical computing and reproducible research.

\subsubsection{Enhancement of Petrophysical Pedagogy}~\\
In petroleum geology education, students often struggle to understand how shifting chromatographic gas fractions relate to complex ratio boundaries and fluid typing. This interactive platform allows students to explore multi-ratio models, customize mathematical expressions, and visualize live log curve responses across standardized wellsite templates in real time.